\documentclass[11pt]{article}

\usepackage[margin=1in]{geometry}
\usepackage{amsmath,amssymb}
\usepackage{booktabs}
\usepackage{hyperref}
\usepackage{xcolor}

\hypersetup{
  colorlinks=true,
  linkcolor=blue!55!black,
  citecolor=blue!55!black,
  urlcolor=blue!55!black
}

\title{Zero-Noise Extrapolation Study for the Split-Basis H$_2$O Matrix Entry}
\author{}
\date{27 July 2026}

\begin{document}
\maketitle

\begin{abstract}
This note summarizes the local noisy-simulator zero-noise extrapolation (ZNE)
study performed for diagonal split-basis matrix entry $(1,1)$ of the H$_2$O
calculation at bond length $1.0$.  It records the simulation and mitigation
pipeline, folding factors, shot-allocation analysis, extrapolation models,
uncertainty treatment, numerical results, and current interpretation.  The
main empirical conclusion is that curved extrapolators substantially reduce
the observed simulator bias for this relatively noisy/deep entry.  The final
five-factor experiment gives the smallest observed RMSE for an
offset-exponential model, closely followed by a quadratic model, while the
linear estimator remains much less variable.  Evidence for curvature is
strong but does not yet meet the conventional $p<0.05$ threshold.
\end{abstract}

\section{Question and setup}

The target is the diagonal matrix element
\[
  H_{11}^{\mathrm{ideal}}=-74.9897034653.
\]
The basis is the orbital-$0$ split basis: the original state $0$ was replaced
by states in fixed orbital sectors, and entry $(1,1)$ is the dominant split
sector studied here.

All experiments reported here were run locally.  Circuits cached by the
sampling workflow were simulated with Qiskit Aer using the noise model and
connectivity of \texttt{FakeQuebec}.  The simulator was configured with the
statevector method, finite shots, and qubit truncation.  This is a noisy,
shot-based simulation, not an ideal statevector expectation-value
calculation.  Readout errors were corrected using independently calibrated
local assignment matrices before the ZNE fit.

The implementation is primarily in
\begin{verbatim}
src/mitigation.py
scripts/run_probabilistic_zne_entry.py
scripts/fit_zne_models.py
\end{verbatim}
and portable results are under
\begin{verbatim}
saved/runtime_runs/
  20260724T024719Z_h2o_1.0_state0_orbital0_split_v1_qwc_fake_quebec/
\end{verbatim}

\section{Noise amplification and execution}

\subsection{Probabilistic local folding}

For an eligible native noisy gate $U$, a local fold replaces it by
\[
  U \longmapsto U U^\dagger U,
\]
which is ideally equivalent to $U$ but executes additional noisy gates.
For fractional amplification $\lambda\in[1,3]$, each eligible gate is folded
with probability
\[
  p_{\mathrm{fold}}=\frac{\lambda-1}{2}.
\]
Different partial folds need not have identical noisy expectation values.
Therefore eight independently randomized folding patterns were used for each
non-unit factor.  The shots assigned to a factor were divided among these
patterns; randomization increased the number of circuits, not the total shots
at that factor.

The final factor grid was
\[
  \lambda\in\{1,\;1.25,\;1.5,\;1.75,\;2\}.
\]
Smaller fractional factors were preferred to full factors $3$ and $5$
because this entry is already noisy and deeper folded circuits can approach
noise saturation.

\paragraph{Current implementation limitation.}
Folding acts on the saved transpiled state-preparation circuits.  Measurement
basis rotations subsequently introduced by the estimator are not presently
folded.  Thus $\lambda$ is an amplification factor for the folded circuit
portion, not a guarantee that every physical error source is multiplied
uniformly.  Realized noisy-gate factors and fold manifests are saved and
should be inspected when interpreting the fits.

\subsection{Readout mitigation}

For each measured physical qubit $q$, calibration circuits prepare
$\lvert0\rangle$ and $\lvert1\rangle$ and estimate the assignment matrix
\[
 A_q =
 \begin{pmatrix}
 p(0|0) & p(0|1)\\
 p(1|0) & p(1|1)
 \end{pmatrix}.
\]
The measured distribution is corrected with $A_q^{-1}$.  For a multi-qubit
Pauli product, the code assumes local assignment errors and applies the
corresponding tensor product of single-qubit inverses.  Readout mitigation was
applied at every noise factor before extrapolation.  It was helpful, but
earlier tests showed that readout error was not the dominant bias source.

\section{Extrapolation and statistical treatment}

\subsection{Models}

Three variance-weighted models were compared:
\begin{align}
 E_{\mathrm{lin}}(\lambda)
   &= \beta_0+\beta_1\lambda,\\
 E_{\mathrm{quad}}(\lambda)
   &= \beta_0+\beta_1\lambda+\beta_2\lambda^2,\\
 E_{\mathrm{exp}}(\lambda)
   &= E_0+s\,\frac{1-\exp(-k\lambda)}{k}.
\end{align}
The last parameterization is an offset exponential written so that
$E_{\mathrm{exp}}(0)=E_0$ and it has a stable linear limit as $k\to0$.
In every case the fitted value at $\lambda=0$ is the ZNE estimate.

If estimates $\widehat E_i$ have variances $v_i$, variance-weighted fitting
minimizes
\[
  \chi^2=\sum_i\frac{
    [\widehat E_i-E_{\mathrm{model}}(\lambda_i)]^2}{v_i}.
\]
For a linear extrapolated estimator
\[
  \widehat E(0)=\sum_i c_i\widehat E_i,
\]
independent factor variances propagate as
\[
  \operatorname{Var}[\widehat E(0)]
  =\sum_i c_i^2v_i.
\]
Factor-level parametric bootstrap resampling was additionally used for all
three fits and is particularly important for the nonlinear exponential fit.
The final results use $10\,000$ bootstrap samples.

\subsection{Shot allocation}

For fixed extrapolation coefficients and single-shot standard deviations
$\sigma_i$, minimizing variance at fixed total shots gives
\[
  N_i \propto |c_i|\sigma_i.
\]
An early equal-allocation run used $147\,317$ shots at each of
$\lambda=\{1,1.5,2\}$.  A pilot-based allocation then assigned
\[
  (N_1,N_{1.5},N_2)
  =(234\,371,\;45\,356,\;162\,224),
\]
while preserving the same total of $441\,951$ shots.  The variance-weighted
linear coefficients after this reallocation became
\[
  (c_1,c_{1.5},c_2)
  =(1.9889,\;0.0222,\;-1.0111).
\]
This illustrates a self-consistency issue: changing $N_i$ changes $v_i$,
which changes a weighted fit and hence changes $c_i$.  Under a strictly
linear model the variance-optimal design tends to place nearly all shots at
the two endpoints; an interior point is nevertheless valuable for detecting
curvature.  A production design should therefore predeclare pilot-derived
weights and reserve a minimum diagnostic allocation at interior factors,
rather than repeatedly adapting weights to the same production data.

\section{Numerical development}

\subsection{Three-factor linear ZNE}

The initial readout-mitigated three-factor run at
$\lambda=\{1,1.5,2\}$ gave
\[
\begin{array}{lccc}
\toprule
 & \text{bias} & \text{variance} & \text{RMSE}\\
\midrule
\lambda=1                         & 2.4949 & 0.00431 & 2.4958\\
\text{equal-allocation linear ZNE}& 1.3022 & 0.02198 & 1.3106\\
\text{optimized weighted ZNE}     & 1.2134 & 0.01546 & 1.2198\\
\bottomrule
\end{array}
\]
The pilot allocation and variance-weighted fit reduced the extrapolated
variance by approximately $29.7\%$ and the RMSE by approximately $6.9\%$
relative to the equal-allocation linear run.

\subsection{Four-factor model comparison}

Adding $\lambda=1.25$ and using $294\,634$ shots per factor produced:
\[
\begin{array}{lrrrr}
\toprule
\text{model} & E(0) & \text{bias} & \text{variance} & \text{RMSE}\\
\midrule
\text{linear}      & -73.7650 & 1.2247 & 0.00943 & 1.2285\\
\text{quadratic}   & -74.3106 & 0.6791 & 0.22591 & 0.8289\\
\text{exponential} & -74.4352 & 0.5545 & 0.52684 & 0.9134\\
\bottomrule
\end{array}
\]
The curved fits reduced observed simulator bias but amplified statistical
variance.  With only four factors, the curvature evidence was not conclusive.

\subsection{Final five-factor independent-seed comparison}

Two independently seeded five-factor runs were inverse-variance pooled.  The
nominal total was $883\,902$ shots per factor.  Resolved folding-pattern
variance was included in the factor variances; therefore nominal shot count
alone was not treated as the effective precision.  The pooled factor data
were
\[
\begin{array}{crr}
\toprule
\lambda & \widehat E(\lambda) & \text{standard deviation}\\
\midrule
1.00 & -72.509300 & 0.02681\\
1.25 & -72.133633 & 0.02773\\
1.50 & -71.774107 & 0.03487\\
1.75 & -71.453188 & 0.02890\\
2.00 & -71.160255 & 0.03185\\
\bottomrule
\end{array}
\]
and the full five-point fits were
\[
\begin{array}{lrrrrr}
\toprule
\text{model} & E(0) & \text{bias} & \text{variance}
             & \text{standard deviation} & \text{RMSE}\\
\midrule
\text{linear}
  & -73.845122 & 1.144581 & 0.00307 & 0.0554 & 1.1459\\
\text{quadratic}
  & -74.321020 & 0.668683 & 0.07574 & 0.2752 & 0.7231\\
\text{offset exponential}
  & -74.405365 & 0.584339 & 0.14510 & 0.3809 & 0.6975\\
\bottomrule
\end{array}
\]
The corresponding fit-quality statistics were
\[
\begin{array}{lrrr}
\toprule
\text{model} & \chi^2/\mathrm{dof} & p\text{-value} & \mathrm{AIC}\\
\midrule
\text{linear}      & 3.129/3 & 0.372 & 7.129\\
\text{quadratic}   & 0.034/2 & 0.983 & 6.034\\
\text{exponential} & 0.051/2 & 0.975 & 6.051\\
\bottomrule
\end{array}
\]
Here
\[
  \mathrm{AIC}=\chi^2+2k
\]
penalizes the residual $\chi^2$ by twice the number $k$ of fit parameters.
Lower AIC is preferred, but differences smaller than roughly $2$ do not
meaningfully distinguish models.  The quadratic has the lowest AIC, but its
advantage over linear is only $1.095$.

\subsection{Held-out prediction and curvature test}

The $\lambda=1.75$ point was excluded from training and then predicted:
\[
\begin{array}{lrrr}
\toprule
\text{model} & \text{prediction} & \text{residual}
             & \text{standardized residual}\\
\midrule
\text{linear}      & -71.475791 & 0.022603 & 0.632\\
\text{quadratic}   & -71.454738 & 0.001550 & 0.041\\
\text{exponential} & -71.456425 & 0.003237 & 0.087\\
\bottomrule
\end{array}
\]
The quadratic predicted the held-out point best, closely followed by the
exponential.  The linear prediction was nevertheless within $0.63$ standard
deviations and therefore was not rejected by this test.

The fitted quadratic curvature was
\[
  \beta_2=-0.2294\pm0.1293,
\]
with bootstrap $95\%$ interval
\[
  [-0.4849,\;0.0254]
\]
and two-sided $p=0.076$.  The evidence for curvature is consequently strong
and reproducible but still falls short of the conventional $p<0.05$
criterion.

\section{What circuit depth implies for model choice}

There is no universal gate-depth boundary at which one should switch from a
linear to a quadratic or exponential extrapolator.  The relevant control
variable is closer to an accumulated circuit error or ``circuit area,'' which
depends jointly on depth, two-qubit error rates, gate composition, observable
support, coherence, and the fidelity with which folding scales those errors.

Majumdar \emph{et al.} varied two-qubit depth and depolarizing error and
constructed calibration phase diagrams for linear, quadratic, exponential,
and ``no-fit'' regions~\cite{majumdar2023}.  Their results support the
following qualitative progression:
\begin{enumerate}
  \item At low depth and low accumulated error, the response is locally
        linear; linear fitting is stable and has the smallest variance.
  \item At greater circuit area, measurable curvature develops and quadratic
        or exponential fits can lower RMSE.
  \item At still larger depth/error, amplified expectation values approach
        noise saturation.  Linear extrapolation can be precisely wrong, while
        exponential fitting can become unstable; no fitted ZNE model is then
        reliable.
\end{enumerate}
For factors $\{1,3,5\}$, the same work found curved extrapolators increasingly
useful at larger depths, including hardware tests out to two-qubit depth
$70$.  With much smaller factors such as $\{1,1.1,1.2\}$, linear fitting
remained useful farther into the deeper/noisier regime and the no-fit region
shrank~\cite{majumdar2023}.  This is consistent with using the present
fractional grid rather than factors $3$ and $5$.

Exponential behavior has a physical motivation under Markovian or Pauli
noise, for which Pauli expectations can decay as one or several
exponentials~\cite{cai2021,krebsbach2022}.  Experiments at large scale have
also reported that exponential extrapolation can outperform linear fitting
when changes among amplified-noise estimates are well resolved, while
downgrading the model when the exponential becomes unstable
\cite{kim2023utility}.  These results do not make exponential fitting
universally correct: coherent, non-Markovian, gate-dependent, leakage, and
nonuniformly amplified errors can violate the assumed response.

For the present entry, the observed flattening, improved curved-model bias,
and held-out prediction are consistent with the ``moderate accumulated
error'' regime in which literature finds curved models useful.  The final
selection is nevertheless not decisive: the exponential has the smallest
observed simulator RMSE, the quadratic has almost the same fit quality with
smaller variance, and the linear fit remains statistically acceptable.

\section{Recommended use}

\begin{enumerate}
  \item Retain all factor estimates and their total variances, including
        folding-pattern variance.
  \item Use predeclared pilot-derived weights; do not select weights or models
        merely because they land closest to a target that would be unknown on
        hardware.
  \item Report the linear estimate as the stable low-variance reference.
        Report quadratic and offset-exponential estimates as curved-model
        alternatives for this deep entry.
  \item Prefer predictive checks, bootstrap intervals, and residual
        diagnostics over choosing a model from circuit depth alone.
  \item Reject ZNE when amplified points have saturated or when fitted
        intercepts are unstable under modest changes of factors/model.
  \item Before hardware deployment, fold the complete physical circuit,
        including measurement-basis rotations, or explicitly quantify the
        unamplified error contribution.
\end{enumerate}

\section{Saved artifacts}

The final independent-seed pooled comparison is stored in
\begin{verbatim}
saved/runtime_runs/
  20260724T024719Z_h2o_1.0_state0_orbital0_split_v1_qwc_fake_quebec/
  probabilistic_zne_entry_1_1_20260727T162129Z/
\end{verbatim}
The directory includes the source run result, calibrated readout data, raw
measurement histograms, folded circuits, numerical ZNE arrays,
\path{zne_model_comparison.json.gz}, and
\path{zne_model_bootstrap_samples.npz}.  A companion independently seeded
high-shot run is
\begin{verbatim}
probabilistic_zne_entry_1_1_20260727T161644Z/
\end{verbatim}

\begin{thebibliography}{9}

\bibitem{giurgica2020}
T.~Giurgica-Tiron, Y.~Hindy, R.~LaRose, A.~Mari, and W.~J.~Zeng,
``Digital zero noise extrapolation for quantum error mitigation,''
\emph{2020 IEEE International Conference on Quantum Computing and
Engineering}, pp.~306--316 (2020).
\href{https://doi.org/10.1109/QCE49297.2020.00045}
{doi:10.1109/QCE49297.2020.00045}.

\bibitem{majumdar2023}
R.~Majumdar, P.~Rivero, F.~Metz, A.~Hasan, and D.~S.~Wang,
``Best practices for quantum error mitigation with digital zero-noise
extrapolation,'' \emph{2023 IEEE International Conference on Quantum
Computing and Engineering}, pp.~881--887 (2023).
\href{https://doi.org/10.1109/QCE57702.2023.00102}
{doi:10.1109/QCE57702.2023.00102}.

\bibitem{cai2021}
Z.~Cai,
``Multi-exponential error extrapolation and combining error mitigation
techniques for NISQ applications,'' \emph{npj Quantum Information}
\textbf{7}, 80 (2021).
\href{https://doi.org/10.1038/s41534-021-00404-3}
{doi:10.1038/s41534-021-00404-3}.

\bibitem{krebsbach2022}
M.~Krebsbach, B.~Trauzettel, and A.~Calzona,
``Optimization of Richardson extrapolation for quantum error mitigation,''
arXiv:2201.08080 (2022).
\url{https://arxiv.org/abs/2201.08080}.

\bibitem{kim2023utility}
Y.~Kim \emph{et al.},
``Evidence for the utility of quantum computing before fault tolerance,''
\emph{Nature} \textbf{618}, 500--505 (2023).
\href{https://doi.org/10.1038/s41586-023-06096-3}
{doi:10.1038/s41586-023-06096-3}.

\end{thebibliography}

\end{document}
